\section{ETL: từ hệ thống nguồn đến kho dữ liệu}\label{sec:etl}
\subsection{ODS và phân tích dữ liệu nguồn}\label{sec:ods}
Dữ liệu từ hai hệ thống nguồn trước hết được chép nguyên trạng vào tầng staging (11 bảng), mỗi dòng gắn mã đợt nạp và thời điểm nạp để truy vết. Có hai chế độ trích xuất: toàn phần (xoá và chép lại) và tăng dần (chỉ lấy dòng mới dựa trên watermark: ngày mua của đơn, ngày liên hệ đầu của MQL, ngày chốt của deal; các bảng con lấy theo đơn mới). Từ staging, các thủ tục SQL làm sạch dữ liệu vào tầng ODS. Bảng toạ độ địa lý 1.000.163 điểm được gom về 19.010 mã bưu chính, lấy median toạ độ và mode tên thành phố sau khi loại 42 điểm nằm ngoài phạm vi Brazil. Tên thành phố được chuẩn hoá chữ thường, bỏ dấu (73.440 điểm toạ độ bị đổi). Cột sai chính tả \code{*\_\allowbreak{}lenght} được đổi tên; 610 sản phẩm không có danh mục được gán \code{unknown}; 13 sản phẩm thuộc hai danh mục thiếu bản dịch được bổ sung bằng bảng seed. Với đơn hàng, thay vì xoá, hệ thống gắn cờ lỗi: 129 trong 45.430 đơn có mốc thời gian sai thứ tự hoặc thiếu ngày (phổ biến nhất là ngày bàn giao vận chuyển trước ngày duyệt, 84 đơn). Các dòng không nạp được đều ghi vào bảng \code{etl\_\allowbreak{}log.rejected\_\allowbreak{}rows} kèm lý do; kiểm thử đối soát cho thấy với mọi bảng số dòng staging bằng số dòng ODS cộng số dòng bị loại.

\begin{table}[htbp]
\centering\small
\caption{Lỗi dữ liệu nguồn được phát hiện khi làm sạch và cách xử lý (số dòng ở trạng thái cuối)}
\label{tab:dq-errors}
\begin{tabular}{p{6.2cm}rp{5.5cm}}
\toprule
\textbf{Lỗi} & \textbf{Số dòng} & \textbf{Xử lý} \\
\midrule
Bàn giao vận chuyển trước khi duyệt thanh toán & 1.359 & Gắn cờ dq\_flags, giữ đơn; khoảng thời gian âm đặt NULL \\
Bàn giao vận chuyển trước ngày mua & 166 & Gắn cờ dq\_flags, giữ đơn; khoảng thời gian âm đặt NULL \\
Giao khách trước khi bàn giao vận chuyển & 23 & Gắn cờ dq\_flags, giữ đơn; khoảng thời gian âm đặt NULL \\
Đơn đã giao nhưng không có ngày duyệt & 14 & Gắn cờ dq\_flags, giữ đơn; khoảng thời gian âm đặt NULL \\
Trạng thái delivered nhưng không có ngày giao & 8 & Gắn cờ dq\_flags, giữ đơn; khoảng thời gian âm đặt NULL \\
Có ngày giao nhưng trạng thái chưa là delivered & 6 & Gắn cờ dq\_flags, giữ đơn; khoảng thời gian âm đặt NULL \\
Đã giao nhưng không có ngày bàn giao & 2 & Gắn cờ dq\_flags, giữ đơn; khoảng thời gian âm đặt NULL \\
Giá trị thanh toán = 0 & 9 & Gắn cờ, chuẩn hoá (số kỳ 0 → 1) \\
Hình thức thanh toán not\_defined & 3 & Gắn cờ, chuẩn hoá (số kỳ 0 → 1) \\
Số kỳ trả góp = 0 (đặt thành 1) & 2 & Gắn cờ, chuẩn hoá (số kỳ 0 → 1) \\
Doanh thu tháng khai báo = 0 & 797 & Gắn cờ, giữ nguyên giá trị \\
Chốt deal trước ngày liên hệ đầu & 1 & Gắn cờ, giữ nguyên giá trị \\
coordinates outside Brazil bounding box & 42 & Loại khỏi ODS, ghi etl\_log.rejected\_rows \\
\bottomrule
\end{tabular}
\\[2pt]{\footnotesize Cờ chất lượng gắn trên đơn, thanh toán và deal; mỗi dòng có thể có nhiều cờ.}
\end{table}

\todo{Nhóm bổ sung phần ``Giới thiệu về ODS'' theo giáo trình (khái niệm, vị trí giữa nguồn và kho, khác biệt với kho) và liên hệ với Chương 1.}

\subsection{Môi trường và mã nguồn}
Hệ thống chạy trên PostgreSQL 18 (cài trực tiếp, không dùng Docker và không dùng PostGIS), Python 3.11 (pandas, SQLAlchemy, psycopg2) cho điều phối và nạp dữ liệu, SQL procedure (PL/pgSQL) cho biến đổi, Power BI Desktop cho dashboard và Git/GitHub để quản lý mã nguồn (nhánh theo từng phase, nhánh \code{main} luôn chạy được). Cấu trúc thư mục: \code{sql/} (DDL, procedure, view, kiểm thử), \code{etl/} (gói Python), \code{scripts/} (xuất dữ liệu, báo cáo), \code{powerbi/} (đặc tả mô hình và DAX), \code{handoff/} và \code{huong-dan/} (tài liệu). Toàn bộ dự án dựng lại từ đầu bằng một lệnh \code{python scripts/rebuild\_all.py --drop} và cho kết quả giống hệt lần trước (Bảng \ref{tab:rebuild}).
\begin{table}[htbp]
\centering\small
\caption{Dựng lại toàn bộ dự án từ đầu bằng một lệnh (kiểm chứng tái lập)}
\label{tab:rebuild}
\begin{tabular}{p{6cm}rp{6cm}}
\toprule
\textbf{Bước} & \textbf{Giây} & \textbf{Kết quả} \\
\midrule
drop olist\_dw + olist\_full & 0.6 & dropped \\
setup\_db (schemas, DDL, procedures, views) & 2.1 & Setup done. \\
load\_source --initial (olist\_dw, to 2017-12-31) & 14.8 & Rejected rows logged in etl\_log.rejected\_rows: 0 (batch 1) \\
load\_source --full (olist\_full) & 23.2 & Rejected rows logged in etl\_log.rejected\_rows: 0 (batch 1) \\
run\_pipeline --mode full (initial state) & 53.3 & tests: 136/136 đạt, 0 lỗi; đối chiếu EDA 6/6 khớp; pipeline 52.4s \\
replay\_increments (11 x simulate + incremental) & 275.3 & 11 đợt hằng tháng (01/2018 → 11/2018); sau đợt cuối: 8,000 MQL, 842 deal, 379 seller có đơn bán được, 99,441 đơn, GMV 13,494,400.74 \\
run\_pipeline --mode incremental (tests + export\_powerbi) & 30.1 & tests: 136/136 đạt, 0 lỗi; đối chiếu EDA 6/6 khớp; pipeline 29.2s \\
\bottomrule
\end{tabular}
\end{table}

\todo{Thêm bảng công nghệ (logo, phiên bản, vai trò), hình cấu trúc thư mục repo, và hướng dẫn cài đặt rút gọn (README.md). Nêu quy ước đặt tên và quy trình làm việc nhóm với Git.}

\subsection{Staging và ODS: làm sạch dữ liệu}\label{sec:staging-ods}
Quy tắc staging và ODS đã mô tả ở mục \ref{sec:ods}; phần này tổng hợp các quy tắc làm sạch và số dòng bị ảnh hưởng.

Các quy tắc làm sạch và số dòng bị ảnh hưởng được tổng hợp ở Bảng \ref{tab:dq-errors} (mục 
ef{sec:ods}). Mỗi dòng bị loại được ghi vào \code{etl\_log.rejected\_rows} cùng lý do; không có dòng nào bị bỏ âm thầm.
\todo{Thêm sơ đồ luồng staging $\rightarrow$ ODS (các job ETL) và bảng ``nguồn $\rightarrow$ đích $\rightarrow$ phép biến đổi'' (Phụ lục C).}

\subsection{Đối soát trước và sau ETL, nạp tăng dần}\label{sec:doi-soat}
Mẫu dữ liệu trước và sau làm sạch được đối chiếu ở các hình dưới đây (nguồn: handoff/before-after-samples.md).
\hinh{3-4-3_before_after_orders.png}{Trước và sau ETL: bảng đơn hàng}{fig:ba-orders}{0.95}
\hinh{3-4-3_before_after_products.png}{Trước và sau ETL: bảng sản phẩm}{fig:ba-products}{0.95}
\hinh{3-4-3_before_after_sellers.png}{Trước và sau ETL: bảng seller}{fig:ba-sellers}{0.95}
\hinh{3-4-3_before_after_geolocation.png}{Trước và sau ETL: bảng toạ độ địa lý}{fig:ba-geo}{0.95}
\hinh{3-4-3_before_after_mql.png}{Trước và sau ETL: bảng MQL}{fig:ba-mql}{0.95}
\hinh{3-4-3_before_after_closed_deals.png}{Trước và sau ETL: bảng deal đã chốt}{fig:ba-deals}{0.95}
\paragraph{Nạp tăng dần.}
Vì bộ dữ liệu Olist là ảnh chụp tĩnh, nhóm mô phỏng việc cập nhật theo tháng. Hệ thống nguồn ban đầu chứa mọi bản ghi đến 31/12/2017 (45.430 đơn); sau đó một chương trình mô phỏng đổ thêm từng tháng vào hệ thống bán hàng và CRM, và kho dữ liệu được nạp tăng dần sau mỗi đợt. Ở mỗi đợt, bước extract chỉ lấy những dòng chưa có ở staging (đơn hàng theo mốc thời gian, các bảng còn lại theo khoá), tầng ODS làm sạch và thử lại các dòng từng bị loại vì thiếu bản ghi cha, chiều khách hàng đóng phiên bản cũ và mở phiên bản mới khi khách đổi mã bưu chính, bảng fact giao dịch thêm dòng, fact snapshot tích luỹ được cập nhật theo mã đơn, còn các snapshot nhỏ được tính lại. Toàn bộ nguồn được đưa vào kho sau 11 đợt hằng tháng (286.144 bản ghi mới), mỗi đợt mất khoảng 17–23 giây gồm cả kiểm thử. Kết quả cuối cùng khớp hoàn toàn với một lần nạp đầy đủ: cả 16 bảng chiều và fact có cùng số dòng, cùng tổng tiền và cùng giá trị kiểm tra nội dung. Mỗi lần chạy đều được ghi vào nhật ký \code{etl\_\allowbreak{}log} (thời gian, số dòng đọc/chèn/cập nhật/loại), làm nguồn cho dashboard giám sát; nhờ đó có thể thấy funnel marketing "lấp đầy" dần: số seller từ funnel có đơn tăng từ 0 lên 379 qua các đợt.

\begin{table}[htbp]
\centering\small
\caption{Các đợt nạp tăng dần: bản ghi mới, thời gian chạy (gồm kiểm thử) và funnel “lấp đầy dần”}
\label{tab:increments}
\begin{tabular}{clrrrrrrrrr}
\toprule
\textbf{Đợt} & \textbf{Đến tháng} & \textbf{Bản ghi mới ở nguồn} & \textbf{ODS chèn} & \textbf{ODS loại} & \textbf{DW chèn / cập nhật} & \textbf{Giây} & \textbf{MQL} & \textbf{Deal} & \textbf{Seller có đơn} & \textbf{Đơn} \\
\midrule
1 & 2018-01 & 38.768 & 38.768 & 0 & 40.535 / 216 & 21,7 & 3.143 & 76 & 3 & 52.699 \\
2 & 2018-02 & 35.979 & 35.979 & 0 & 37.817 / 255 & 21,3 & 4.171 & 189 & 30 & 59.427 \\
3 & 2018-03 & 38.659 & 38.659 & 0 & 40.668 / 311 & 22,2 & 5.345 & 336 & 55 & 66.638 \\
4 & 2018-04 & 37.515 & 37.515 & 0 & 39.700 / 489 & 23,5 & 6.697 & 543 & 117 & 73.577 \\
5 & 2018-05 & 37.079 & 37.079 & 0 & 39.402 / 450 & 24,4 & 8.000 & 665 & 199 & 80.450 \\
6 & 2018-06 & 32.035 & 32.035 & 0 & 34.594 / 468 & 25,3 & 8.000 & 722 & 281 & 86.617 \\
7 & 2018-07 & 32.489 & 32.489 & 0 & 35.245 / 435 & 24,8 & 8.000 & 759 & 341 & 92.909 \\
8 & 2018-08 & 33.490 & 33.490 & 0 & 36.399 / 369 & 26,2 & 8.000 & 792 & 380 & 99.421 \\
9 & 2018-09 & 87 & 87 & 0 & 3.195 / 23 & 21,1 & 8.000 & 815 & 380 & 99.437 \\
10 & 2018-10 & 37 & 37 & 0 & 54 / 21 & 20,7 & 8.000 & 836 & 380 & 99.441 \\
11 & 2018-11 & 6 & 6 & 0 & 18 / 6 & 20,7 & 8.000 & 842 & 380 & 99.441 \\
\bottomrule
\end{tabular}
\end{table}

\paragraph{Thử lại dòng mồ côi.} Dòng con bị loại vì chưa có dòng cha (ví dụ deal đến trước MQL) được thử lại ở các đợt sau và biến mất khỏi danh sách loại khi đã nạp. \textbf{Lưu ý trung thực}: cơ chế này mới được chứng minh bằng kịch bản kiểm thử; trong 11 đợt thật không có dòng mồ côi nào xảy ra vì nguồn luôn nạp dòng cha trước.
\todo{Thêm sơ đồ data pipeline hoặc ETL jobs (có logo công nghệ); ví dụ trước/sau cho thêm bảng nếu giảng viên yêu cầu; kết quả so sánh incremental với full: 16/16 bảng khớp (handoff/2026-11-05-P5.md).}
